\documentclass[11pt,a4paper]{article}

% PACKAGES
\usepackage[utf8]{inputenc}
\usepackage[T1]{fontenc}
\usepackage[slovak]{babel}
\usepackage{amsmath, amssymb}
\usepackage{amsthm} % For theorem environments
\usepackage{graphicx}
\usepackage{enumerate}
\usepackage{hyperref}

\usepackage{tikz}
\usetikzlibrary{tikzmark}
\usetikzlibrary{arrows.meta}
\usetikzlibrary{calc}
\usetikzlibrary{shapes.geometric} % Library required for ellipse shape
\usetikzlibrary{positioning}      % Library for easy relative positioning
\usetikzlibrary{cd}
\usetikzlibrary{babel}
\usetikzlibrary{decorations.pathreplacing}

\usepackage[all,pdf,2cell]{xy}
\usepackage{soul}
\usepackage{framed}
\usepackage{array}
\usepackage{todonotes}
\usepackage{nicematrix}

\NiceMatrixOptions
  { pgf-node-code = \pgfsetfillcolor{white} \pgfusepathqfill }

\pgfset{nicematrix/cell-node/.append style = { inner sep = 1pt } }

% THEOREM ENVIRONMENTS SETUP
% Style for definitions and examples (upright text, bold title)
\theoremstyle{definition} 
\newtheorem{definicia}{Definícia}[section]
\newtheorem{priklad}[definicia]{Príklad}

% Style for theorems (italicized text, bold title)
\theoremstyle{plain} 
\newtheorem{tvrdenie}[definicia]{Tvrdenie}
\newtheorem{veta}[definicia]{Veta}
\newtheorem{dosledok}[definicia]{Dôsledok}

% Style for remarks and questions (upright text, italic title)
\theoremstyle{remark} 
\newtheorem*{question}{Otázka}
\newtheorem*{pozn}{Poznámka}

% Macros

% Simple abbreviations
\newcommand{\id}{\mathrm{id}}
\newcommand{\Nat}{\mathbb{N}}
\newcommand{\Set}{\mathbf{Set}}
\newcommand{\R}{\mathbb{R}}
\newcommand{\ora}[1]{\overrightarrow{#1}}
\newcommand{\Lo}{\mathrm{Lo}} % Notation for Linearny Obal

% More complex stuff
% Restriction of a function (from stackexchange)
\newcommand\restr[2]{{% we make the whole thing an ordinary symbol
  \left.\kern-\nulldelimiterspace % automatically 
resize the bar with \right
  #1 % the function
  \littletaller % pretend it's a little taller at normal size
  \right|_{#2} % this is the delimiter
  }}
\newcommand{\littletaller}{\mathchoice{\vphantom{\big|}}{}{}{}}
\setcounter{section}{-1} 
% DOCUMENT START
\begin{document}
\title{Lineárna algebra 1\\(texty k prednáškam)}
\author{Gejza Jenča}
\date{Verzia 1}
\maketitle
\tableofcontents

\section{Polia}

Do tejto chvíle sme pod vektorovým priestorom rozumeli takúto štvoricu $(V, +, \cdot, \vec{0})$, kde:
\begin{itemize}
    \item $V$ je množina (prvky sú vektory),
    \item $+: V \times V \to V$ je sčítanie,
    \item $\cdot: \R \times V \to V$ je násobenie skalárom,
    \item $\vec{0} \in V$ je nulový vektor.
\end{itemize}

To nestačilo, aby to bol vektorový priestor. Museli o týchto veciach platiť aj nejaké vlastnosti (axiómy). Pomocou týchto axióm sme potom začali budovať teóriu vektorových priestorov.

Keďže sme pri budovaní teórie používali iba axiómy, všetky pojmy a tvrdenia o nich budú fungovať pre všetky vektorové priestory, ktoré spĺňajú axiómy.

Teraz ideme zaviesť do celej veci ďalšiu flexibilitu a toto rozhýbeme.

\begin{tvrdenie}[Pozorovanie]
Úplne všetky veci, ktoré sme dokazovali v prvom semestri by fungovali, ak by sme vymenili $\R$ za niečo, čoho prvky vieme sčítať, odčítať, násobiť! To sme nepochybne so skalármi robili.
\end{tvrdenie}

Čo ešte je nejaká množina, ktorej prvky vieme sčítať, odčítať, násobiť?
\begin{itemize}
    \item Odpoveď: $\Nat$? Mohli by sme robiť lineárnu algebru s takýmito skalármi? Nie. Potrebujeme deliť (napríklad keď riešime sústavu lineárnych rovníc).
    \item Naproti tomu $\R$ má ešte celú kopu vlastností, ktoré sme nepoužili, aspoň kladnosť čísel, konvergencia postupností a iné veci.
\end{itemize}

\begin{question}
Aké vlastnosti sčítania, násobenia v $\R$ sme pre budovanie našej teórie potrebovali?
\end{question}
Odpoveď je táto definícia.

\begin{definicia}
\textbf{Pole} je pätica $(K, +, \cdot, 0, 1)$, kde:
\begin{itemize}
    \item $K$ je množina,
    \item $+: K \times K \to K$ a $\cdot: K \times K \to K$ sú binárne operácie,
    \item $0 \in K$ a $1 \in K$ sú konštanty,
\end{itemize}
pričom sú splnené nasledujúce axiómy poľa:
\begin{enumerate}
    \item $\forall a,b,c \in K: (a+b)+c = a+(b+c)$ (asociatívnosť sčítania)
    \item $\forall a,b \in K: a+b = b+a$ (komutatívnosť sčítania)
    \item $\forall a \in K: a+0 = a$ (existencia neutrálneho prvku pre sčítanie)
    \item $\forall a \in K, \exists b \in K: a+b=0$ (toto $b$ je nutne jediné a značí sa $-a$, existencia opačného prvku)
    \item $\forall a,b,c \in K: (a \cdot b) \cdot c = a \cdot (b \cdot c)$ (asociatívnosť násobenia)
    \item $\forall a \in K: 1 \cdot a = a$ (existencia neutrálneho prvku pre násobenie)
    \item $\forall a \in K, a \neq 0, \exists b \in K: a \cdot b = 1$ (toto $b$ je nutne jediné a značí sa $a^{-1}$, existencia inverzného prvku)
    \item $\forall a,b,c \in K: a \cdot (b+c) = (a \cdot b) + (a \cdot c)$ (distributívnosť)
    \item $0 \neq 1$ (netriviálnosť)
\end{enumerate}
\end{definicia}

\subsection*{Príklady polí}
Dva príklady sú nám dobre známe: reálne čísla $\R$ a racionálne čísla $\mathbb{Q}$. Väčšinu času tejto prednášky strávime zavedením a skúmaním ďalšieho príkladu: komplexných čísel $\mathbb{C}$.

Predtým ako sa do toho pustíme, spomeniem ešte jednu triedu príkladov: konečné polia $Z_p$.

Nech $p$ je prvočíslo. Vybavme množinu $Z_p = \{0, 1, \dots, p-1\}$ operáciami sčítania a násobenia takto:
\begin{itemize}
    \item $a \oplus b = (a+b) \pmod p$
    \item $a \odot b = (a \cdot b) \pmod p$
\end{itemize}

Napríklad pre $p=3$ dostaneme:

\begin{center}
\begin{tabular}{c|ccc}
$\oplus$ & 0 & 1 & 2 \\
\hline
0 & 0 & 1 & 2 \\
1 & 1 & 2 & 0 \\
2 & 2 & 0 & 1
\end{tabular}
\hspace{2cm}
\begin{tabular}{c|ccc}
$\odot$ & 0 & 1 & 2 \\
\hline
0 & 0 & 0 & 0 \\
1 & 0 & 1 & 2 \\
2 & 0 & 2 & 1
\end{tabular}
\end{center}

Potom $(Z_p, \oplus, \odot, 0, 1)$ je pole (nedokazujem). Platí, že ak $p$ nie je
prvočíslo, axiómy poľa splnené nie sú.  Polia $Z_p$ sú veľmi užitočné objekty (teória
kódovania, kryptografia, \dots). My sa nimi však príliš zaoberať nebudeme.

\section{Komplexné čísla}

Cardano v roku 1545 vzorec pre riešenie kubickej rovnice obsahuje záporné odmocniny
aj keď sú všetky riešenia reálne. Výpočet vzorca teda predpokladal existenciu
odmocniny zo záporného čísla.

Vezmime reálne čísla $\R$. Pridajme nový prvok $i \notin \R$ a zadefinujme $i^2 =
-1$. Pridávajme nové prvky tak, aby sme pridali len to, čo je nutné a pritom dostali
pole. Pribudnú postupne $bi$, $a+bi$ kde $a, b \in \R$

\begin{definicia}[Komplexné čísla]
$\mathbb{C} = \R \times \R$ (ako množina). Prvky ale nezapisujeme $(a, b)$, ale takto:
$$ a + bi $$
kde $a, b \in \R$. Toto sa nazýva \textbf{algebrický tvar}.
V skutočnosti:
\begin{itemize}
    \item $a \equiv (a, 0)$
    \item $i \equiv (0, 1)$
\end{itemize}
Sčítanie $+: \mathbb{C} \times \mathbb{C} \to \mathbb{C}$ je definované:
$$ (a_1 + b_1 i) + (a_2 + b_2 i) = (a_1 + a_2) + (b_1 + b_2)i $$
Násobenie $\cdot: \mathbb{C} \times \mathbb{C} \to \mathbb{C}$ je definované:
$$ (a_1 + b_1 i)(a_2 + b_2 i) = (a_1 a_2 - b_1 b_2) + (a_1 b_2 + b_1 a_2)i $$
(Toto je predpis pre násobenie dvojíc).
\end{definicia}

Nulový prvok je $0 = 0 + 0i$. Jednotkový prvok je $1 = 1 + 0i$.
$( \mathbb{C}, +, \cdot, 0, 1 )$ je potom pole.

\paragraph{Náznak dôkazu:} Axiómy pre $+$ sú jasné. Pre násobenie: ak $z = a+bi$ a $z \neq 0$, ako vyzerá $z^{-1}$?
$$ z^{-1} = \frac{1}{z} = \frac{1}{a+bi} $$
Rozšírme zlomok výrazom $a-bi$:
$$ \frac{1}{a+bi} \cdot \frac{a-bi}{a-bi} = \frac{a-bi}{a^2 - abi + abi + b^2} = \frac{a-bi}{a^2+b^2} = \frac{a}{a^2+b^2} - \frac{b}{a^2+b^2}i $$
Teda
$$ (a+bi)^{-1} = \frac{a}{a^2+b^2} - \frac{b}{a^2+b^2}i $$
Treba overiť, že to vyjde: $(a+bi) \cdot z^{-1} = 1$. (DÚ).

Iné horeuvedené operácie sčítania a násobenia komplexných čísel je možné jednoducho (ale trochu prácne) overiť. Nebudeme to robiť.

\subsection*{Geometrická interpretácia}
Keďže komplexné číslo $a+bi$ je stotožniteľné s dvojicou čísel $(a,b)$, ponúka sa možnosť ho reprezentovať ako bod v rovine vybavenej dvoma kolmými súradnicovými osami:
\begin{itemize}
    \item Reálna časť $z$ ($\mathrm{Re}\, z$) na osi x.
    \item Imaginárna časť $z$ ($\mathrm{Im}\, z$) na osi y.
\end{itemize}

Sčítanie komplexných čísel potom dostáva zrejmý geometrický význam analogicky ako sčítanie geometrických vektorov v rovine.
Násobenie komplexných čísel je ale tiež interpretovateľné geometricky.

\begin{definicia}
\begin{itemize}
    \item \textbf{Absolútna hodnota} $z = a+bi$ je $|z| = \sqrt{a^2+b^2}$.
    \item \textbf{Konjugované číslo} (združené) k $z$ je $\overline{z} = a-bi = \mathrm{Re}\, z - (\mathrm{Im}\, z)i$.
\end{itemize}
\end{definicia}

Platí $z \cdot \overline{z} = (a+bi)(a-bi) = a^2+b^2 \in \R$. Teda $|z|^2 = z \cdot \overline{z}$ alebo $|z| = \sqrt{z \cdot \overline{z}}$.
Všimnime si, že hore je napísané, že $z^{-1} = \frac{\overline{z}}{|z|^2}$.

\subsection*{Goniometrický tvar}
Nahliadnutím do pravouhlého trojuholníka v rovine vidíme, že:
$$ \mathrm{Im}\, z = |z| \sin(\arg z) $$
$$ \mathrm{Re}\, z = |z| \cos(\arg z) $$
kde $\arg z$ je uhol, ktorý zviera sprievodič bodu $z$ s kladnou polosou x.
Z toho dostaneme, že $z = \mathrm{Re}\, z + i \mathrm{Im}\, z = |z|\cos(\arg z) + i|z|\sin(\arg z)$.
Teda
$$ z = |z|(\cos \varphi + i \sin \varphi) $$
kde $\varphi = \arg z$. Toto sa nazýva \textbf{goniometrický tvar} komplexného čísla.

Vypočítajme, aký vplyv majú polárne tvary na násobenie komplexných čísel. Nech $z = |z|(\cos \alpha + i \sin \alpha)$ a $w = |w|(\cos \beta + i \sin \beta)$.
$$ z \cdot w = |z||w| (\cos \alpha + i \sin \alpha)(\cos \beta + i \sin \beta) $$
$$ = |z||w| [ (\cos \alpha \cos \beta - \sin \alpha \sin \beta) + i (\sin \alpha \cos \beta + \cos \alpha \sin \beta) ] $$
Použitím súčtových vzorcov:
$$ \cos(\alpha+\beta) = \cos \alpha \cos \beta - \sin \alpha \sin \beta $$
$$ \sin(\alpha+\beta) = \sin \alpha \cos \beta + \cos \alpha \sin \beta $$
dostaneme:
$$ z \cdot w = |z||w| (\cos(\alpha+\beta) + i \sin(\alpha+\beta)) $$

Tento vzorec značí geometricky, že pri násobení \textbf{sčítame argumenty} a \textbf{vynásobíme absolútne hodnoty}.
$$ |zw| = |z||w| $$
$$ \arg(zw) = \arg z + \arg w $$

\begin{priklad}
Nech $z = 1+i$, $w = \sqrt{3}+i$.
$|z| = \sqrt{2}$, $\arg z = \frac{\pi}{4}$.
$|w| = 2$, $\arg w = \frac{\pi}{6}$.
Potom $z \cdot w$ má absolútnu hodnotu $2\sqrt{2}$ a argument $\frac{\pi}{4} + \frac{\pi}{6} = \frac{5\pi}{12}$.
\end{priklad}

\subsection*{Moivreova veta}
Rozvinutím vzorca pre násobenie dostaneme vzorec pre mocninu komplexného čísla. Pre $n \in \Nat$:
$$ z^n = |z|^n (\cos(n \arg z) + i \sin(n \arg z)) $$

\subsection*{Riešenie rovnice $z^n = c$}
Predpokladajme teraz, že máme $n \in \Nat$ a číslo $c \in \mathbb{C}$. Máme za úlohu riešiť rovnicu $z^n = c$ (t.j. hľadať $n$-tú odmocninu z $c$).
Upravme obe strany na goniometrický tvar. Nech $z = |z|(\cos \varphi + i \sin \varphi)$ a $c = |c|(\cos \psi + i \sin \psi)$.
Potom rovnica $z^n = c$ má tvar:
$$ |z|^n (\cos n\varphi + i \sin n\varphi) = |c|(\cos \psi + i \sin \psi) $$
Dostaneme $|z|^n = |c|$, teda $|z| = \sqrt[n]{|c|}$.
Pre uhly platí:
$$ \cos n\varphi = \cos \psi $$
$$ \sin n\varphi = \sin \psi $$
Aké sú možné hodnoty $n\varphi$?
Jedným z riešení je $\psi$, ale čo ostatné? Funkcie sin a cos majú periódu $2\pi$. Z toho dostaneme:
$$ n\varphi = \psi + 2k\pi, \quad k \in \{0, 1, \dots, n-1\} $$
Teda
$$ \varphi_k = \frac{\psi + 2k\pi}{n} $$

Geometricky riešenia tvoria pravidelný $n$-uholník, ležia na kružnici s polomerom $\sqrt[n]{|c|}$.

\begin{priklad}
Riešme rovnicu $z^3 = -2 + 2i$.
Máme $c = -2+2i$. $|c| = \sqrt{4+4} = \sqrt{8}$.
Argument $\arg c = \frac{3\pi}{4}$ (lebo reálna časť záporná, imaginárna kladná, rovnaká veľkosť).
Riešenia teda ležia na kružnici so stredom v 0 a polomerom $\sqrt[3]{\sqrt{8}} = \sqrt{2}$.
Argumenty riešení sú:
$$ \varphi_k = \frac{\frac{3\pi}{4} + 2k\pi}{3} = \frac{\pi}{4} + \frac{2k\pi}{3} $$
Pre $k=0: \varphi_0 = \frac{\pi}{4}$. Riešenie $z_0 = \sqrt{2}(\cos \frac{\pi}{4} + i \sin \frac{\pi}{4}) = \sqrt{2}(\frac{\sqrt{2}}{2} + i\frac{\sqrt{2}}{2}) = 1+i$.
Pre $k=1, 2$ dostaneme ostatné vrcholy trojuholníka.
\end{priklad}

\end{document}
